% =============================================================================
%  Rapport TP2 DevSecOps - Intégration de la sécurité dans un pipeline CI/CD
%  Compilation : pdflatex rapport.tex (2 fois, pour la table des matières)
%                ou importer le dossier « rapport » dans Overleaf.
%  Captures   : déposer les images PNG dans rapport/captures/ avec le nom
%               indiqué dans chaque cadre gris. Tant qu'une image manque,
%               un cadre de remplacement décrivant la capture est affiché.
% =============================================================================
\documentclass[11pt,a4paper]{article}

\usepackage[utf8]{inputenc}
\usepackage[T1]{fontenc}
\usepackage[french]{babel}
\usepackage{lmodern}
\usepackage{microtype}
\usepackage{amssymb}
\usepackage[margin=2.3cm]{geometry}
\usepackage{graphicx}
\usepackage[table]{xcolor}
\usepackage{listings}
\usepackage{booktabs}
\usepackage{tabularx}
\usepackage{array}
\usepackage{float}
\usepackage{caption}
\usepackage{enumitem}
\usepackage{fancyhdr}
\usepackage{tcolorbox}
\tcbuselibrary{skins,breakable}
\usepackage{tikz}
\usetikzlibrary{arrows.meta,positioning,shapes.geometric,fit,backgrounds,babel}
\usepackage[hidelinks]{hyperref}

% ---------------------------------------------------------------- Informations
\newcommand{\etudiant}{A. CHABBAH}          % <-- à adapter (Prénom NOM)
\newcommand{\formation}{DevSecOps}          % <-- à adapter (formation / promo)
\newcommand*{\enseignant}{}                 % <-- à adapter (nom de l'enseignant)
\newcommand{\depot}{https://github.com/amin119/TP2-DEVSECOPS}

% ---------------------------------------------------------------- Couleurs
\definecolor{bleu}{HTML}{1F4E79}
\definecolor{vert}{HTML}{2E7D32}
\definecolor{rouge}{HTML}{C62828}
\definecolor{orange}{HTML}{EF6C00}
\definecolor{gris}{HTML}{F3F4F6}
\definecolor{grisfonce}{HTML}{6B7280}
\definecolor{codebg}{HTML}{F8F8F8}

% ---------------------------------------------------------------- En-têtes
\pagestyle{fancy}
\fancyhf{}
\lhead{\small TP2 -- Sécurité dans un pipeline CI/CD}
\rhead{\small \etudiant}
\cfoot{\thepage}
\setlength{\headheight}{14pt}

% ---------------------------------------------------------------- Listings
\lstdefinelanguage{yaml}{
  keywords={true,false,null,on},
  sensitive=false,
  comment=[l]{\#},
  morestring=[b]",
  morestring=[b]',
}
\lstset{
  basicstyle=\ttfamily\footnotesize,
  backgroundcolor=\color{codebg},
  frame=single, rulecolor=\color{black!15},
  keywordstyle=\color{bleu}\bfseries,
  commentstyle=\color{grisfonce}\itshape,
  stringstyle=\color{vert},
  numbers=left, numberstyle=\tiny\color{grisfonce}, numbersep=6pt,
  breaklines=true, breakatwhitespace=false, columns=fullflexible,
  showstringspaces=false, tabsize=2, xleftmargin=1.5em, framexleftmargin=1.5em,
  literate=
    {é}{{\'e}}1 {è}{{\`e}}1 {ê}{{\^e}}1 {ë}{{\"e}}1 {à}{{\`a}}1 {â}{{\^a}}1
    {î}{{\^i}}1 {ï}{{\"i}}1 {ô}{{\^o}}1 {ù}{{\`u}}1 {û}{{\^u}}1 {ç}{{\c c}}1
    {É}{{\'E}}1 {È}{{\`E}}1 {À}{{\`A}}1 {Ç}{{\c C}}1
    {«}{{\guillemotleft}}1 {»}{{\guillemotright}}1 {’}{{'}}1 {–}{{--}}1 {—}{{---}}1
}

% ---------------------------------------------------------------- Captures
% \capture{nom-fichier}{légende}{ce qu'il faut capturer}
\newcommand{\capture}[3]{%
  \begin{figure}[H]
    \centering
    \IfFileExists{captures/#1.png}{%
      \includegraphics[width=\linewidth,height=0.42\textheight,keepaspectratio]{captures/#1.png}%
    }{%
      \begin{tcolorbox}[width=\linewidth, colback=gris, colframe=grisfonce,
                        arc=2mm, boxrule=0.6pt, halign=center, valign=center, height=5.2cm]
        {\Large\color{grisfonce} Capture d'écran à insérer}\\[2mm]
        \texttt{captures/#1.png}\\[3mm]
        \begin{minipage}{0.85\linewidth}\centering\small #3\end{minipage}
      \end{tcolorbox}%
    }
    \caption{#2}
    \label{fig:#1}
  \end{figure}
}

% ---------------------------------------------------------------- Encadrés
\newtcolorbox{objectif}{colback=bleu!5, colframe=bleu, arc=1mm, boxrule=0.8pt,
  left=3mm, right=3mm, fonttitle=\bfseries, title=Objectif de l'étape}
\newtcolorbox{regle}{colback=rouge!4, colframe=rouge, arc=1mm, boxrule=0.8pt,
  left=3mm, right=3mm, fonttitle=\bfseries, title=Règle bloquante}
\newtcolorbox{bilan}{colback=vert!5, colframe=vert, arc=1mm, boxrule=0.8pt,
  left=3mm, right=3mm, fonttitle=\bfseries, title=Bilan}

\newcommand{\ok}{\textcolor{vert}{\textbf{Succès}}}
\newcommand{\ko}{\textcolor{rouge}{\textbf{Échec}}}
\newcommand{\fichier}[1]{\texttt{#1}}

\begin{document}

% =============================================================================
%  PAGE DE GARDE
% =============================================================================
\begin{titlepage}
  \centering
  \vspace*{1.5cm}
  {\large \formation\par}
  \vspace{2.5cm}
  {\color{bleu}\rule{\linewidth}{1.2pt}\par}
  \vspace{0.6cm}
  {\Huge\bfseries Intégration de la sécurité\\[3mm] dans un pipeline CI/CD\par}
  \vspace{0.4cm}
  {\Large Travaux pratiques n\textsuperscript{o}\,2 -- DevSecOps\par}
  \vspace{0.6cm}
  {\color{bleu}\rule{\linewidth}{1.2pt}\par}
  \vspace{1.5cm}
  {\large SAST $\cdot$ SCA $\cdot$ Secrets $\cdot$ DAST $\cdot$ Security gate\par}
  \vspace{0.4cm}
  {\normalsize GitHub Actions -- Semgrep -- Bandit -- Gitleaks -- Trivy -- OWASP ZAP\par}
  \vfill
  \begin{tabular}{rl}
    \textbf{Réalisé par :} & \etudiant \\
    \ifx\enseignant\empty\else \textbf{Encadré par :} & \enseignant \\ \fi
    \textbf{Dépôt :} & \href{\depot}{\texttt{\detokenize{github.com/amin119/TP2-DEVSECOPS}}} \\
    \textbf{Date :} & \today \\
  \end{tabular}
  \vspace{1cm}
\end{titlepage}

\tableofcontents
\newpage
\listoffigures
\newpage

% =============================================================================
\section{Introduction}
% =============================================================================

L'approche \textbf{DevSecOps} consiste à intégrer la sécurité à chaque étape du cycle de
développement (« \emph{shift left} ») plutôt qu'à la traiter en fin de projet. Le pipeline
d'intégration et de déploiement continus (CI/CD) devient alors un \textbf{garde-fou} : chaque
modification du code est automatiquement compilée, testée, puis analysée par plusieurs outils de
sécurité complémentaires. Si une vulnérabilité critique ou de haute sévérité est détectée, la
modification ne peut pas être fusionnée et l'application n'est pas déployée.

Ce TP a pour objectifs :
\begin{itemize}[itemsep=2pt]
  \item d'automatiser les tests de sécurité (SAST, SCA, DAST) dans un pipeline CI/CD ;
  \item de comprendre le rôle du pipeline comme barrière de sécurité ;
  \item de configurer des règles de validation qui bloquent les déploiements non sécurisés ;
  \item de démontrer le fonctionnement de ces règles par une \emph{pull request} volontairement
        vulnérable, puis par sa correction.
\end{itemize}

\subsection{Application support : SecureNotes}

Pour ce TP, nous avons développé \textbf{SecureNotes}, une petite application web
\textbf{Python/Flask} de prise de notes, avec :
\begin{itemize}[itemsep=2pt]
  \item une interface HTML (liste, création et consultation de notes) protégée par jeton CSRF ;
  \item une API JSON (\fichier{GET/POST /api/notes}, \fichier{GET /api/notes/<id>}) ;
  \item une route de supervision \fichier{/health} ;
  \item une base SQLite accédée exclusivement par requêtes paramétrées ;
  \item des en-têtes de sécurité HTTP (CSP, \fichier{X-Frame-Options}, \fichier{X-Content-Type-Options},
        \fichier{Referrer-Policy}, \fichier{Permissions-Policy}, COOP/CORP/COEP\ldots) ;
  \item une image Docker durcie (utilisateur non privilégié, mises à jour système, \emph{healthcheck}).
\end{itemize}

\begin{table}[H]
\centering
\caption{Structure du dépôt}
\begin{tabularx}{\linewidth}{>{\ttfamily}l X}
\toprule
\textrm{\textbf{Chemin}} & \textbf{Rôle} \\
\midrule
app/                       & Code de l'application Flask (routes, accès BD, sécurité, gabarits) \\
tests/                     & Tests unitaires \texttt{pytest} \\
Dockerfile                 & Image d'exécution (Gunicorn, utilisateur non-root) \\
requirements*.txt          & Dépendances figées (production / développement) \\
.github/workflows/ci.yml   & Pipeline DevSecOps GitHub Actions \\
scripts/semgrep\_gate.py   & Porte de sécurité SAST (analyse du rapport Semgrep) \\
scripts/zap\_gate.py       & Porte de sécurité DAST (analyse du rapport ZAP) \\
.zap/rules.tsv             & Configuration des règles ZAP (alertes informatives ignorées) \\
rapport/                   & Le présent rapport \\
\bottomrule
\end{tabularx}
\end{table}

\capture{cap01-depot-github}{Dépôt GitHub du projet}{Page d'accueil du dépôt GitHub
(\texttt{github.com/amin119/TP2-DEVSECOPS}) montrant l'arborescence des fichiers
et le README.}

% =============================================================================
\section{Architecture du pipeline DevSecOps}
% =============================================================================

Le pipeline est défini dans \fichier{.github/workflows/ci.yml}. Il se déclenche sur chaque
\emph{push} vers \fichier{main} et sur chaque \emph{pull request} ciblant \fichier{main}.
La figure~\ref{fig:pipeline} en présente l'enchaînement.

\begin{figure}[H]
\centering
\begin{tikzpicture}[
  node distance=6mm and 9mm,
  job/.style={draw, rounded corners=2pt, minimum width=3.3cm, minimum height=1.05cm,
              align=center, font=\small, fill=white, thick},
  base/.style={job, draw=bleu, fill=bleu!8},
  secu/.style={job, draw=orange, fill=orange!8},
  gate/.style={job, draw=rouge, fill=rouge!8, minimum width=3cm},
  dep/.style={job, draw=vert, fill=vert!8, minimum width=3cm},
  fl/.style={-{Stealth[length=2.5mm]}, thick, color=black!60}
]
  \node[base] (build) {\textbf{Build \& tests}\\flake8, pytest};
  \node[secu, right=14mm of build, yshift=24mm] (sast) {\textbf{SAST}\\Semgrep, Bandit};
  \node[secu, below=of sast] (secrets) {\textbf{Secrets}\\Gitleaks};
  \node[secu, below=of secrets] (sca) {\textbf{SCA}\\Trivy, pip-audit};
  \node[secu, below=of sca] (dast) {\textbf{DAST}\\OWASP ZAP};
  \node[gate, right=14mm of build, xshift=40mm] (gate) {\textbf{Security}\\\textbf{gate}};
  \node[dep, right=of gate] (deploy) {\textbf{Déploiement}\\image GHCR};

  \foreach \j in {sast,secrets,sca,dast} {
    \draw[fl] (build.east) -- ++(5mm,0) |- (\j.west);
    \draw[fl] (\j.east) -- ++(5mm,0) |- (gate.west);
  }
  \draw[fl] (gate) -- node[above, font=\scriptsize]{si succès} node[below, font=\scriptsize]{\& main} (deploy);

  \node[font=\scriptsize, color=grisfonce, below=1mm of dast, align=center]
    {environnement éphémère\\(conteneur Docker sur le runner)};
\end{tikzpicture}
\caption{Enchaînement des jobs du pipeline DevSecOps}
\label{fig:pipeline}
\end{figure}

Les quatre contrôles de sécurité s'exécutent \textbf{en parallèle} après la réussite du build : le
développeur reçoit ainsi un retour complet en une seule exécution (tous les problèmes sont remontés
en même temps, et non un par un). Un job d'agrégation, la \textbf{security gate}, échoue si l'un des
contrôles a échoué ; c'est ce seul job qui est déclaré obligatoire dans la protection de la branche
\fichier{main}. Le déploiement n'a lieu que sur \fichier{main}, après une gate réussie.

\begin{table}[H]
\centering
\caption{Outils utilisés et règles bloquantes}
\label{tab:outils}
\small
\begin{tabularx}{\linewidth}{l l X X}
\toprule
\textbf{Type} & \textbf{Outil} & \textbf{Périmètre analysé} & \textbf{Condition d'échec} \\
\midrule
Qualité & flake8, pytest & code Python, tests unitaires & erreur de lint ou test en échec \\
SAST & Semgrep & code source (règles \texttt{p/python}, \texttt{p/flask}, \texttt{p/owasp-top-ten}) & résultat de sévérité \texttt{ERROR}/\texttt{HIGH}/\texttt{CRITICAL} \\
SAST & Bandit & code source Python & problème de sévérité \texttt{HIGH} (confiance $\geq$ \texttt{MEDIUM}) \\
Secrets & Gitleaks & fichiers du dépôt & au moins un secret détecté \\
SCA & Trivy (fs) & dépendances (\texttt{requirements*.txt}) & CVE \texttt{HIGH} ou \texttt{CRITICAL} \\
SCA & Trivy (image) & image Docker (OS + paquets) & CVE \texttt{HIGH}/\texttt{CRITICAL} disposant d'un correctif \\
SCA & pip-audit & dépendances (base OSV) & informatif, non bloquant \\
DAST & OWASP ZAP & application en cours d'exécution & alerte de risque \texttt{High} \\
\bottomrule
\end{tabularx}
\end{table}

\paragraph{Remarque sur les niveaux de sévérité.} OWASP ZAP n'a pas de niveau « critique » : son
échelle est \emph{Informational}, \emph{Low}, \emph{Medium}, \emph{High}. Le niveau \emph{High}
(le plus élevé) est donc assimilé à « critique ». Le seuil est paramétrable par la variable
\fichier{ZAP\_FAIL\_RISK} (3 = High, 2 = Medium).

% =============================================================================
\section{Étape 1 -- Pipeline de base (build et tests unitaires)}
% =============================================================================

\begin{objectif}
Disposer d'un pipeline CI qui installe les dépendances, vérifie la qualité du code et exécute les
tests unitaires à chaque modification.
\end{objectif}

Le premier job, \fichier{build-test}, contient les étapes suivantes :
\begin{enumerate}[itemsep=2pt]
  \item \textbf{checkout} du dépôt (\fichier{actions/checkout}) ;
  \item installation de Python 3.12 avec cache pip (\fichier{actions/setup-python}) ;
  \item installation des dépendances figées (\fichier{requirements-dev.txt}) ;
  \item analyse de style avec \textbf{flake8} ;
  \item exécution des \textbf{tests unitaires} avec \textbf{pytest} et mesure de couverture
        (\fichier{pytest-cov}), le rapport \fichier{coverage.xml} étant archivé comme artefact.
\end{enumerate}

\begin{lstlisting}[language=yaml, caption={Job \texttt{build-test} (extrait de \texttt{ci.yml})}]
build-test:
  name: Build & tests unitaires
  runs-on: ubuntu-latest
  steps:
    - uses: actions/checkout@v5
    - uses: actions/setup-python@v6
      with:
        python-version: ${{ env.PYTHON_VERSION }}
        cache: pip
    - name: Installation des dépendances
      run: pip install -r requirements-dev.txt
    - name: Lint (flake8)
      run: flake8 app tests wsgi.py
    - name: Tests unitaires + couverture
      run: pytest -v --cov=app --cov-report=term-missing --cov-report=xml
\end{lstlisting}

Les tests unitaires (9 tests sur la version de base) vérifient le fonctionnement de l'application
mais aussi certaines propriétés de sécurité : présence des en-têtes HTTP, rejet des formulaires sans
jeton CSRF, échappement HTML du contenu des notes, validation des entrées de l'API.
La couverture obtenue est de \textbf{98\,\%}.

\paragraph{Éléments notables du fichier de configuration.}
\begin{itemize}[itemsep=2pt]
  \item \fichier{permissions: contents: read} au niveau global : principe du moindre privilège pour
        le jeton \fichier{GITHUB\_TOKEN} ; seuls les jobs qui en ont besoin obtiennent des droits
        supplémentaires (\fichier{security-events: write}, \fichier{packages: write}).
  \item \fichier{concurrency} : une nouvelle exécution sur la même branche annule la précédente.
  \item Versions de dépendances \textbf{figées} : les builds sont reproductibles et l'analyse SCA
        porte exactement sur les versions déployées.
\end{itemize}

\capture{cap02-pipeline-base}{Exécution réussie du pipeline de base}{Onglet \emph{Actions} :
exécution du commit « Étape 1 » avec le job \emph{Build \& tests unitaires} en vert.}

\capture{cap03-tests-unitaires}{Journal des tests unitaires et de la couverture}{Journal du job
\emph{Build \& tests unitaires}, étape « Tests unitaires + couverture » dépliée : liste des tests
\texttt{PASSED} et tableau de couverture.}

% =============================================================================
\section{Étape 2 -- Intégration du SAST et du SCA}
% =============================================================================

\begin{objectif}
Automatiser l'analyse statique du code (SAST), la recherche de secrets et l'analyse de la composition
logicielle (SCA), et bloquer le pipeline en cas de vulnérabilité critique ou haute.
\end{objectif}

\subsection{SAST : Semgrep et Bandit}

Deux outils complémentaires sont utilisés :
\begin{itemize}[itemsep=2pt]
  \item \textbf{Semgrep}, moteur multi-langages basé sur des règles, avec les jeux de règles
        \fichier{p/python}, \fichier{p/flask} et \fichier{p/owasp-top-ten} (analyse de
        flux de données « \emph{taint} » pour les injections) ;
  \item \textbf{Bandit}, analyseur spécialisé pour Python (appels dangereux, \fichier{shell=True},
        désérialisation, cryptographie faible\ldots).
\end{itemize}

Les rapports sont produits au format \textbf{SARIF} et publiés dans l'onglet \emph{Security
$\rightarrow$ Code scanning} de GitHub, ce qui permet de visualiser les alertes directement sur les
lignes de code concernées. Ils sont également archivés comme artefacts.

\begin{regle}
\begin{itemize}[leftmargin=*, itemsep=1pt]
  \item Semgrep : échec si au moins un résultat a la sévérité \texttt{ERROR}, \texttt{HIGH} ou
        \texttt{CRITICAL} (script \fichier{scripts/semgrep\_gate.py}).
  \item Bandit : échec si au moins un problème de sévérité \texttt{HIGH} avec une confiance
        $\geq$ \texttt{MEDIUM} : \fichier{bandit -r app --severity-level high --confidence-level medium}.
\end{itemize}
\end{regle}

Les portes sont implémentées en deux temps : les étapes de « gate » sont marquées
\fichier{continue-on-error} afin que \emph{tous} les rapports soient produits et publiés, puis une
étape finale « Verdict » fait échouer le job si l'une des portes a échoué.

\begin{lstlisting}[language=yaml, caption={Portes de sécurité du job SAST}]
- name: "Security gate SAST : Semgrep (ERROR/HIGH/CRITICAL)"
  id: semgrep_gate
  continue-on-error: true
  run: python scripts/semgrep_gate.py reports/semgrep.json

- name: "Security gate SAST : Bandit (sévérité HIGH)"
  id: bandit_gate
  continue-on-error: true
  run: bandit -r app wsgi.py --severity-level high --confidence-level medium

- name: Verdict SAST
  run: |
    if [ "${{ steps.semgrep_gate.outcome }}" != "success" ] || \
       [ "${{ steps.bandit_gate.outcome }}" != "success" ]; then
      echo "::error::Vulnérabilités SAST critiques/hautes détectées - fusion bloquée."
      exit 1
    fi
\end{lstlisting}

\subsection{Détection de secrets : Gitleaks}

\textbf{Gitleaks} recherche les secrets (clés d'API, jetons, mots de passe, clés privées) présents
dans les fichiers du dépôt. Toute détection fait échouer le job \fichier{secrets}. Dans l'application,
la clé secrète Flask n'est jamais codée en dur : elle est lue depuis la variable d'environnement
\fichier{SECRET\_KEY} ou générée aléatoirement au démarrage.

\subsection{SCA : Trivy et pip-audit}

L'analyse de composition logicielle vérifie que les composants tiers ne contiennent pas de
vulnérabilités connues (CVE) :
\begin{itemize}[itemsep=2pt]
  \item \textbf{Trivy -- mode \fichier{fs}} : analyse des fichiers de dépendances Python ;
  \item \textbf{Trivy -- mode \fichier{image}} : analyse de l'image Docker construite (paquets du
        système Debian \emph{et} paquets Python installés) ;
  \item \textbf{pip-audit} : second avis basé sur la base OSV/PyPI, à titre informatif (il ne
        permet pas de filtrer par sévérité).
\end{itemize}

\begin{regle}
Trivy échoue (\fichier{--exit-code 1}) si une CVE de sévérité \texttt{HIGH} ou \texttt{CRITICAL}
est trouvée dans les dépendances. Pour l'image, l'option \fichier{--ignore-unfixed} limite le blocage
aux vulnérabilités pour lesquelles un correctif existe (une vulnérabilité sans correctif ne peut pas
être résolue par l'équipe et bloquerait indéfiniment le pipeline).
\end{regle}

\begin{lstlisting}[language=yaml, caption={Portes de sécurité du job SCA}]
- name: "Security gate SCA : dépendances (HIGH, CRITICAL)"
  id: fs_gate
  continue-on-error: true
  run: $TRIVY fs --scanners vuln --severity HIGH,CRITICAL --exit-code 1 /src

- name: "Security gate SCA : image (HIGH, CRITICAL corrigeables)"
  id: image_gate
  continue-on-error: true
  run: |
    $TRIVY image --scanners vuln --ignore-unfixed --severity HIGH,CRITICAL \
      --exit-code 1 securenotes:${{ github.sha }}
\end{lstlisting}

\capture{cap04-pipeline-sast-sca}{Pipeline avec les jobs SAST, secrets et SCA}{Graphe de
l'exécution du commit « Étape 2 » : \emph{Build \& tests} suivi des jobs \emph{SAST},
\emph{Secrets} et \emph{SCA}, tous en vert.}

\capture{cap05-sast-resume}{Résumé SAST sur la branche principale}{Page \emph{Summary} de
l'exécution : tableau « SAST -- Semgrep » (0 résultat bloquant) ; éventuellement le journal
Bandit « No issues identified ».}

\capture{cap06-sca-trivy}{Résultat de l'analyse Trivy sur la branche principale}{Journal du job
\emph{SCA}, étapes « Trivy -- rapport complet des dépendances » et « rapport de l'image » :
tableaux sans vulnérabilité HIGH/CRITICAL.}

% =============================================================================
\section{Étape 3 -- Analyse dynamique (DAST) sur environnement éphémère}
% =============================================================================

\begin{objectif}
Déployer l'application dans un environnement éphémère, l'attaquer avec OWASP ZAP pendant son
exécution, et faire échouer le pipeline si des failles critiques sont découvertes.
\end{objectif}

\subsection{Environnement éphémère}

Le job \fichier{dast} construit l'image Docker de l'application puis la démarre sur le
\emph{runner} GitHub Actions avec une clé secrète aléatoire générée à la volée
(\fichier{openssl rand -hex 32}). Le pipeline attend que la route \fichier{/health} réponde avant de
lancer le scan. Cet environnement est \textbf{éphémère} : il n'existe que pendant la durée du job, il
est isolé de la production, et il est détruit explicitement à la fin (étape exécutée même en cas
d'échec grâce à \fichier{if: always()}), puis avec la machine virtuelle du runner.

\subsection{Scan OWASP ZAP}

Nous utilisons le \textbf{full scan} de ZAP (\fichier{zap-full-scan.py}) plutôt que le
\emph{baseline scan} : le baseline est uniquement passif (il observe les réponses), alors que le
full scan enchaîne un \emph{spider} (exploration des liens et formulaires) puis un \textbf{scan
actif} qui envoie de véritables charges d'attaque (injection SQL, XSS, injection de commandes,
\emph{path traversal}\ldots). Seul le scan actif permet de détecter les failles d'injection.

\begin{lstlisting}[language=yaml, caption={Déploiement éphémère et scan ZAP}]
- name: Déploiement de l'environnement éphémère
  run: |
    docker run -d --name securenotes-ephemeral -p 8000:8000 \
      -e SECRET_KEY="$(openssl rand -hex 32)" securenotes:dast
    for i in $(seq 1 30); do
      if curl -fsS "$TARGET/health"; then exit 0; fi
      sleep 2
    done
    exit 1

- name: Scan OWASP ZAP (full scan - spider + attaque active)
  run: |
    docker run --rm --network host -v "$PWD/zap:/zap/wrk:rw" \
      ghcr.io/zaproxy/zaproxy:stable \
      zap-full-scan.py -t "$TARGET" -c rules.tsv -m 2 -T 15 -I \
      -r zap-report.html -J zap-report.json -w zap-report.md

- name: "Security gate DAST : alertes de risque High"
  run: python3 scripts/zap_gate.py zap/zap-report.json
\end{lstlisting}

Options principales : \fichier{-m 2} limite le spider à 2 minutes, \fichier{-T 15} borne la durée
totale du scan, \fichier{-c rules.tsv} ignore trois règles purement informatives, et \fichier{-I}
empêche ZAP d'échouer sur de simples avertissements : la décision est déléguée au script
\fichier{zap\_gate.py}.

\begin{regle}
Le script \fichier{scripts/zap\_gate.py} lit le rapport JSON de ZAP, affiche un tableau de toutes les
alertes dans le résumé de l'exécution, et termine en erreur si au moins une alerte a un niveau de
risque \texttt{High} (seuil configurable par \fichier{ZAP\_FAIL\_RISK}).
\end{regle}

Les rapports HTML, JSON et Markdown de ZAP sont archivés comme artefacts du job.

\capture{cap07-pipeline-complet}{Pipeline complet réussi sur la branche principale}{Graphe de
l'exécution du commit « Étape 3 » : \emph{Build}, \emph{SAST}, \emph{Secrets}, \emph{SCA},
\emph{DAST}, \emph{Security gate} et \emph{Déploiement}, tous en vert.}

\capture{cap08-dast-resume}{Résumé DAST sur la version saine}{Page \emph{Summary} : tableau
« DAST -- OWASP ZAP » avec 0 alerte High (seules des alertes Low/Informational éventuelles).}

\capture{cap09-zap-rapport-html}{Rapport HTML d'OWASP ZAP}{Télécharger l'artefact
\texttt{zap-reports}, ouvrir \texttt{zap-report.html} dans le navigateur et capturer le récapitulatif
des alertes par niveau de risque.}

% =============================================================================
\section{Security gate, protection de branche et déploiement}
% =============================================================================

\subsection{Agrégation des contrôles}

Le job \fichier{security-gate} dépend de tous les contrôles (\fichier{needs}) et s'exécute toujours
(\fichier{if: always()}), même lorsqu'un contrôle a échoué. Il publie un tableau récapitulatif et
échoue si un seul résultat est différent de \fichier{success}.

\subsection{Protection de la branche \texttt{main}}

Pour que le pipeline soit réellement un garde-fou, la branche \fichier{main} est protégée :
\begin{itemize}[itemsep=2pt]
  \item fusion uniquement via \emph{pull request} ;
  \item \textbf{status check obligatoire} : \fichier{Security gate} ;
  \item la branche doit être à jour avant fusion.
\end{itemize}
Le bouton « \emph{Merge} » d'une pull request reste donc grisé tant que la security gate est en
échec.

\capture{cap10-protection-branche}{Règle de protection de la branche main}{\emph{Settings
$\rightarrow$ Branches} (ou \emph{Rules $\rightarrow$ Rulesets}) : règle sur \texttt{main} avec
« Require status checks to pass » et le check \emph{Security gate} sélectionné.}

\subsection{Déploiement conditionnel}

Le job \fichier{deploy} ne s'exécute que sur un \emph{push} vers \fichier{main} et uniquement si la
security gate a réussi. Il construit l'image validée et la publie dans le registre GitHub Container
Registry (\fichier{ghcr.io}), étiquetée avec le SHA du commit. Une version non sûre ne peut donc
jamais être publiée.

% =============================================================================
\section{Étape 4 -- Validation : pull request vulnérable puis corrigée}
% =============================================================================

\begin{objectif}
Démontrer que le pipeline bloque une modification vulnérable, puis qu'il accepte la version
corrigée.
\end{objectif}

\subsection{Pull request n\textsuperscript{o}\,1 : fonctionnalité vulnérable}

La branche \fichier{feature/recherche-vulnerable} ajoute une fonctionnalité de recherche de notes,
une route de diagnostic réseau et un import de notes au format YAML. Plusieurs vulnérabilités y ont
été introduites \textbf{volontairement} :

\begin{lstlisting}[language=Python, caption={Extraits de code vulnérable (branche \texttt{feature/recherche-vulnerable})}]
@bp.get("/search")
def search():
    q = request.args.get("q", "")
    # VULNERABLE (CWE-89) : requête SQL construite par concaténation
    rows = conn.execute(
        f"SELECT id, title, content FROM notes WHERE title LIKE '%{q}%' ..."
    ).fetchall()
    # VULNERABLE (CWE-79) : saisie réinjectée sans échappement
    return render_template("search.html", query=Markup(q), results=Markup(items))

@bp.get("/admin/ping")              # VULNERABLE (CWE-306) : aucune authentification
def admin_ping():
    host = request.args.get("host", "127.0.0.1")
    # VULNERABLE (CWE-78) : injection de commande système
    result = subprocess.run(f"ping -c 1 {host}", shell=True, ...)

@bp.post("/api/notes/import")
def api_import_notes():
    # VULNERABLE (CWE-502) : désérialisation YAML non sûre
    data = yaml.load(request.get_data(as_text=True), Loader=yaml.Loader)
\end{lstlisting}

De plus, la version de \textbf{Gunicorn} est rétrogradée à \fichier{21.2.0}, affectée par les
vulnérabilités \textbf{CVE-2024-1135} et \textbf{CVE-2024-6827} (\emph{HTTP request smuggling},
sévérité \texttt{HIGH}).

\paragraph{Point important.} Les tests unitaires de cette branche \textbf{passent} (la recherche et
l'import fonctionnent) : sans les contrôles de sécurité, cette modification aurait été fusionnée et
déployée.

\begin{table}[H]
\centering
\caption{Vulnérabilités introduites et contrôles censés les détecter}
\label{tab:matrice}
\small
\begin{tabularx}{\linewidth}{X l c c c}
\toprule
\textbf{Vulnérabilité} & \textbf{CWE} & \textbf{SAST} & \textbf{SCA} & \textbf{DAST} \\
\midrule
Injection SQL dans \texttt{/search}                & CWE-89  & \checkmark & & \checkmark \\
XSS réfléchi dans \texttt{/search}                 & CWE-79  & \checkmark & & \checkmark \\
Route admin sans authentification + injection de commande (\texttt{/admin/ping}) & CWE-306 / CWE-78 & \checkmark & & \checkmark \\
Désérialisation YAML non sûre                       & CWE-502 & \checkmark & & \\
Divulgation d'erreurs SQL                           & CWE-209 & & & \checkmark \\
Gunicorn 21.2.0 (CVE-2024-1135, CVE-2024-6827)      & CWE-444 & & \checkmark & \\
\bottomrule
\end{tabularx}
\end{table}

Ce tableau illustre la \textbf{complémentarité} des approches : le SAST voit le code mais pas
l'environnement d'exécution, le SCA voit les composants tiers mais pas le code métier, et le DAST
confirme l'exploitabilité réelle depuis l'extérieur, sans connaître le code.

\capture{cap11-pr-vulnerable-checks}{Pull request vulnérable : checks en échec et fusion bloquée}{Page
de la PR n\textsuperscript{o}\,1, en bas : liste des checks avec SAST, SCA, DAST et \emph{Security gate}
en rouge, et le message « Merging is blocked » / bouton de fusion désactivé.}

\capture{cap12-pr-vulnerable-graphe}{Graphe d'exécution de la pull request vulnérable}{Onglet
\emph{Checks} ou \emph{Actions} de la PR n\textsuperscript{o}\,1 : graphe du workflow montrant le build
vert et les jobs de sécurité en échec.}

\capture{cap13-sast-echec}{Détections SAST sur la pull request vulnérable}{Journal du job
\emph{SAST} : résultat de Bandit (B602 \texttt{subprocess\_popen\_with\_shell\_equals\_true},
sévérité High) et tableau Semgrep des règles déclenchées dans le résumé.}

\capture{cap14-sca-echec}{Détection SCA : Gunicorn vulnérable}{Journal du job \emph{SCA}, étape
« Security gate SCA : dépendances » : tableau Trivy listant \texttt{gunicorn 21.2.0} avec
CVE-2024-1135 / CVE-2024-6827 (HIGH).}

\capture{cap15-dast-echec}{Détections DAST : alertes High remontées par ZAP}{Résumé de l'exécution,
section « DAST -- OWASP ZAP » : alertes High (injection SQL, XSS réfléchi, injection de commande\ldots)
et message « pipeline bloqué ».}

\capture{cap16-code-scanning}{Alertes dans GitHub Code Scanning}{Onglet \emph{Security
$\rightarrow$ Code scanning} (ou annotations dans l'onglet \emph{Files changed} de la PR) : alertes
Semgrep/Bandit positionnées sur les lignes vulnérables.}

\capture{cap17-security-gate-echec}{Security gate en échec}{Résumé de l'exécution, tableau
« Security gate » : Build \texttt{success}, SAST/SCA/DAST \texttt{failure}.}

\subsection{Pull request n\textsuperscript{o}\,2 : version corrigée}

La branche \fichier{fix/recherche-securisee} propose la même fonctionnalité, écrite de manière sûre :

\begin{table}[H]
\centering
\caption{Corrections apportées}
\small
\begin{tabularx}{\linewidth}{l X}
\toprule
\textbf{Vulnérabilité} & \textbf{Correction} \\
\midrule
Injection SQL     & requête paramétrée (\texttt{LIKE ?}) dans \texttt{db.search\_notes()} \\
XSS réfléchi      & suppression de \texttt{Markup()} : échappement automatique de Jinja2 \\
Route \texttt{/admin/ping} & route supprimée (fonctionnalité non nécessaire, aucune commande système) \\
YAML non sûr      & \texttt{yaml.safe\_load()} + validation du schéma (liste de 1 à 50 notes, champs contrôlés) \\
Erreurs SQL       & plus d'exception SQL possible ; aucune trace technique renvoyée au client \\
Gunicorn          & retour à la version maintenue \texttt{26.2.0} \\
Saisie            & longueur de la recherche limitée à 100 caractères \\
\bottomrule
\end{tabularx}
\end{table}

\begin{lstlisting}[language=Python, caption={Code corrigé (branche \texttt{fix/recherche-securisee})}]
def search_notes(term):                       # app/db.py
    pattern = f"%{term}%"
    return get_db().execute(
        "SELECT id, title, content FROM notes WHERE title LIKE ? OR content LIKE ?",
        (pattern, pattern),
    ).fetchall()

@bp.get("/search")                            # app/routes.py
def search():
    q = request.args.get("q", "").strip()[:MAX_QUERY]
    results = db.search_notes(q) if q else []
    return render_template("search.html", query=q, results=results)

@bp.post("/api/notes/import")
def api_import_notes():
    try:
        data = yaml.safe_load(request.get_data(as_text=True))
    except yaml.YAMLError:
        return jsonify(error="YAML invalide"), 400
    ...  # validation de chaque note avec _validate()
\end{lstlisting}

Des \textbf{tests de non-régression de sécurité} ont été ajoutés : la charge
\texttt{x' OR '1'='1' --} ne renvoie aucune note, la charge \texttt{<script>} est échappée, la route
\texttt{/admin/ping} renvoie 404, et un document YAML contenant
\texttt{!!python/object/apply:os.system} est rejeté (15 tests, 96\,\% de couverture).

\capture{cap18-pr-corrigee-checks}{Pull request corrigée : tous les checks passent}{Page de la PR
n\textsuperscript{o}\,2 : tous les checks en vert et le message « All checks have passed », bouton
« Merge pull request » actif.}

\capture{cap19-pr-corrigee-dast}{Résumé de sécurité de la pull request corrigée}{Résumé de
l'exécution de la PR n\textsuperscript{o}\,2 : tableaux SAST (0 bloquant), DAST (0 High) et
Security gate entièrement à \texttt{success}.}

\capture{cap20-deploiement}{Fusion et déploiement de la version corrigée}{Après fusion de la PR
n\textsuperscript{o}\,2 : exécution sur \texttt{main} avec le job \emph{Déploiement} en vert, et/ou
la page \emph{Packages} du dépôt montrant l'image publiée sur \texttt{ghcr.io}.}

% =============================================================================
\section{Synthèse des résultats}
% =============================================================================

\begin{table}[H]
\centering
\caption{Résultats des exécutions du pipeline}
\label{tab:resultats}
\small
\begin{tabular}{l c c c c c c c}
\toprule
\textbf{Exécution} & \textbf{Build} & \textbf{SAST} & \textbf{Secrets} & \textbf{SCA} & \textbf{DAST} & \textbf{Gate} & \textbf{Déploiement} \\
\midrule
\texttt{main} -- étape 3          & \ok & \ok & \ok & \ok & \ok & \ok & \ok \\
PR n\textsuperscript{o}\,1 (vulnérable) & \ok & \ko & \ok & \ko & \ko & \ko & -- \\
PR n\textsuperscript{o}\,2 (corrigée)   & \ok & \ok & \ok & \ok & \ok & \ok & -- \\
\texttt{main} après fusion        & \ok & \ok & \ok & \ok & \ok & \ok & \ok \\
\bottomrule
\end{tabular}
\end{table}

\begin{bilan}
Le pipeline remplit son rôle de garde-fou : la pull request vulnérable, dont les tests unitaires
passent pourtant, est bloquée par trois contrôles indépendants (SAST, SCA et DAST) et ne peut pas
être fusionnée. La version corrigée franchit tous les contrôles, est fusionnée, puis déployée
automatiquement.
\end{bilan}

% =============================================================================
\section{Analyse critique et pistes d'amélioration}
% =============================================================================

\subsection{Limites constatées}
\begin{itemize}[itemsep=2pt]
  \item \textbf{Faux positifs / faux négatifs} : les outils SAST reposent sur des règles ; certaines
        détections peuvent être non pertinentes (ex. Bandit B108 sur un chemin temporaire) et
        nécessitent un tri. Inversement, une faille logique (contrôle d'accès métier) peut échapper
        aux trois outils.
  \item \textbf{Durée du DAST} : le scan actif est le job le plus long ; il est borné
        (\texttt{-m 2 -T 15}) pour rester compatible avec un cycle de PR. Un scan plus profond
        (authentifié, avec définition OpenAPI) pourrait être planifié la nuit.
  \item \textbf{Vulnérabilités sans correctif} dans l'image de base : elles sont volontairement
        ignorées par la gate (\texttt{--ignore-unfixed}) mais restent visibles dans les rapports.
  \item \textbf{Seuils} : bloquer uniquement sur \emph{High/Critical} est un compromis entre sécurité
        et vélocité ; les alertes \emph{Medium} doivent être suivies dans un backlog.
\end{itemize}

\subsection{Améliorations possibles}
\begin{itemize}[itemsep=2pt]
  \item Génération d'un \textbf{SBOM} (CycloneDX/SPDX) et \textbf{signature de l'image} (Cosign)
        pour garantir l'intégrité de la chaîne d'approvisionnement ;
  \item analyse de l'\textbf{infrastructure as code} (Dockerfile, workflows) avec Trivy
        \texttt{config}, Checkov ou Hadolint ;
  \item épinglage des actions GitHub par \textbf{SHA} de commit et Dependabot pour leurs mises à jour ;
  \item centralisation des résultats dans \textbf{DefectDojo} ou SonarQube pour le suivi dans le temps ;
  \item scan DAST \textbf{authentifié} et piloté par la spécification OpenAPI de l'API ;
  \item fichier d'exceptions versionné (avec justification et date d'expiration) pour gérer les
        faux positifs sans désactiver les contrôles.
\end{itemize}

% =============================================================================
\section{Conclusion}
% =============================================================================

Ce TP a permis de construire un pipeline CI/CD complet dans lequel la sécurité est un critère
d'acceptation au même titre que les tests fonctionnels. L'association du SAST (Semgrep, Bandit), de
la détection de secrets (Gitleaks), du SCA (Trivy) et du DAST (OWASP ZAP sur un environnement
éphémère) offre une défense en profondeur : chaque outil couvre des angles morts des autres. Les
règles bloquantes, la security gate et la protection de la branche \texttt{main} transforment ces
analyses en un véritable garde-fou : une vulnérabilité critique ne peut plus atteindre la
production sans être explicitement corrigée.

% =============================================================================
\appendix
\section{Annexe -- Fichier complet du pipeline}
% =============================================================================
\lstinputlisting[language=yaml, caption={\texttt{.github/workflows/ci.yml}}]{annexes/ci.yml}

\section{Annexe -- Scripts de security gate}
\lstinputlisting[language=Python, caption={\texttt{scripts/zap\_gate.py}}]{annexes/zap_gate.py}
\lstinputlisting[language=Python, caption={\texttt{scripts/semgrep\_gate.py}}]{annexes/semgrep_gate.py}

\end{document}
